\section{Histograms and binning}

A histogram is a graphical summary of data. It is especially useful for
continuous or nearly continuous observations, where exact repeated values may be
rare and a table of exact values may not show the shape of the data clearly.

The basic idea is to divide the range of the data into intervals, called bins,
and count how many observations fall into each bin.

\subsection*{Bins and counts}

Suppose the bins are
\[
[a_0,a_1),\quad [a_1,a_2),\quad \ldots,\quad [a_{m-1},a_m).
\]
The count in the \(j\)-th bin is
\[
n_j=\#\{i:a_{j-1}\leq x_i<a_j\}.
\]
The relative frequency in that bin is
\[
\frac{n_j}{n}.
\]
This is the proportion of observations that fall in the bin.

For example, if \(n=100\) and \(18\) observations fall in the interval
\([2,4)\), then the relative frequency of that bin is
\[
\frac{18}{100}=0.18.
\]
This means that \(18\%\) of the observed data lie in that interval.

\subsection*{Frequency histograms and relative-frequency histograms}

A frequency histogram uses bar heights equal to bin counts \(n_j\). A
relative-frequency histogram uses bar heights equal to \(n_j/n\). Both are useful,
but they answer slightly different visual questions.

Frequency histograms show how many observations are in each bin. Relative-
frequency histograms show what proportion of the sample is in each bin. When
comparing samples of different sizes, relative frequencies are usually easier to
compare than raw counts.

\subsection*{Density-scaled histograms}

When a histogram is compared with a probability density, the bar heights should
usually be scaled so that bar areas represent relative frequencies. If the
\(j\)-th bin has width
\[
\Delta_j=a_j-a_{j-1},
\]
then the density-scaled height is
\[
h_j=\frac{n_j}{n\Delta_j}.
\]
The area of the bar is then
\[
h_j\Delta_j=\frac{n_j}{n}.
\]
Thus the area of the bar equals the empirical probability of that bin.

\begin{remarkbox}
A density-scaled histogram should be read by areas, just as a probability density
is read by areas. The height of a bar depends on the bin width; the area of a bar
represents the proportion of observations in that bin.
\end{remarkbox}

\subsection*{Histogram versus density}

A histogram is computed from data. A density is part of a theoretical probability
model. They may be compared, but they are not the same object.

For example, if a sample is generated from a normal distribution, its histogram
may look bell-shaped. The normal density is a smooth theoretical curve. The
histogram is a finite-sample visual summary. If another sample is generated from
the same model, the histogram will usually look slightly different.

This variability is not an error. It is the natural effect of finite sampling.

\begin{center}
\begin{tabular}{lll}
\toprule
Object & Built from & Meaning \\
\midrule
Histogram & observed data & empirical visual summary \\
Density & probability model & theoretical distribution shape \\
\bottomrule
\end{tabular}
\end{center}

\subsection*{The choice of bins matters}

A histogram depends on the choice of bins. Wide bins produce a smoother but less
detailed picture. Narrow bins show more detail but may look irregular, especially
for small samples.

Therefore a histogram is not a purely mechanical truth about a dataset. It is a
summary that depends on a modelling choice: how the data are grouped. Good bin
choices make important structure visible without creating misleading patterns.

\subsection*{A computer experiment}

A useful way to understand histograms is to generate samples from a known model.
For example, one may generate \(1000\) values from a standard normal distribution
and draw a density-scaled histogram. The histogram will usually resemble the
normal density, but it will not be identical to it.

A typical Python-style workflow is:
\begin{verbatim}
import numpy as np

sample = np.random.normal(loc=0, scale=1, size=1000)
sample_mean = np.mean(sample)
sample_std = np.std(sample)
\end{verbatim}

The computer experiment does not prove a theorem. It creates a controlled
example where the theoretical model is known, so that the relationship between a
finite sample and the underlying distribution can be observed.

\begin{summarybox}
A histogram groups data into bins. A density-scaled histogram is designed so that
bar areas equal relative frequencies. Histograms are empirical and depend on bin
choice; densities are theoretical objects from probability models.
\end{summarybox}
